\section{从基础 T2V 到编辑、个性化与音频}

基础模型学到的视觉规律可以支持多种应用，但不同任务需要不同输入接口与训练数据。本章用六个代表状态区分：纯文本生成关注时空 plausibility；编辑关注内容保持；个性化关注身份；音频关注跨模态同步。把它们都称为“Movie Gen”容易掩盖模型家族内部的训练差异。

\subsection{模型是否真的学到了“视觉世界”}

前面的架构与数据章节解释了模型怎样训练，本节回到输出，检验它究竟学到了什么。讲者用两段样例说明规模化模型能够生成相机运动、物体运动和某些物理关系。读这些图时应寻找连续性与相互作用，而不是只看单帧纹理；同时要保留证据边界：观察数据可以产生强规律拟合，不自动产生因果可干预的世界模型。

\lecturefigure{slide-257-learned-visual-world.jpg}{风筝奔跑样例：主体、风筝、海滩和相机运动共同演化}{00:41:17}

该片段要求模型保持人物身份与衣着，让风筝与运动方向协调，并处理镜头跟随。它展示了多因素联合生成，但不能回答风力、绳张力等隐变量是否被显式表示。

\lecturefigure{slide-305-physics-sample.jpg}{泳池样例：水面、甜甜圈、树懒与饮料形成接触和浮力关系}{00:42:11}

视觉上合理的接触关系说明模型学到大量统计常识。更严格的物理评估应测试反事实：改变物体质量、液体、碰撞速度或约束条件，观察结果是否稳定遵循规律，而非只在常见场景中看起来合理。

\begin{warningbox}{“World model”至少有三种强度}
弱定义是能生成视觉上合理的未来；中等定义要求对干预和反事实保持一致；强定义还要求可用于规划和预测。Movie Gen 样例支持第一层并触及第二层现象，但不能仅凭演示声称具备完整因果模拟器。
\end{warningbox}

\teachervoice{00:41:11--00:42:52，老师说模型似乎从观看视频中学到对象、运动和物理规律；他把这视为规模化的惊喜，而不是对显式物理推理的证明。}

\subsection{精确编辑：生成质量之外还要守住原视频}

纯文本生成只需构造一个合理世界，编辑还必须尊重已经存在的世界。本节因此把评价拆成“改对了什么”和“误改了什么”：编辑模型同时接收原视频与文本指令，必须把变化限制在目标区域和属性。下面两组四宫格展示局部替换、删除、增加和风格变换；它们比单一输出更适合检查哪些内容被保留。

\lecturefigure{slide-337-precise-editing-vr.jpg}{VR 场景编辑：替换、移除头显并增加视觉效果}{00:42:54}

同一输入支持多种指令，说明编辑接口能够解耦 source content 与 edit intent。评估应分为 edit fidelity 和 preservation：前者看目标变化，后者看人物、动作、背景和镜头是否无意漂移。

\lecturefigure{slide-374-precise-editing-penguin.jpg}{企鹅视频编辑：服装、环境物体与铅笔画风格}{00:43:54}

跨全局风格和局部对象的编辑难度不同。增加伞是局部结构修改，穿服装要求与运动同步，铅笔风格则重绘全画面但应保持场景布局。一个统一分数会隐藏这些失败模式。

\begin{knowledgebox}{编辑数据为什么稀缺}
真实世界几乎没有“同一视频编辑前/编辑后 + 精确指令”的天然 pairs。课堂提示 Movie Gen 使用自监督构造来产生训练信号。没有这种数据设计，基础 T2V 模型即使生成很强，也难以精确保持输入内容。
\end{knowledgebox}

\subsection{个性化与音频属于不同条件路径}

编辑任务之后，本节继续比较两条完全不同的条件扩展：个性化视频增加 reference identity，音频生成增加 video/text-to-audio 条件。两者都扩展基础生成能力，却面对不同目标：个性化要平衡身份与动作，音频要对齐视觉事件、节奏和环境语义；因此不能用同一“视频质量”分数概括。

\lecturefigure{slide-388-personalization.jpg}{个性化视频：参考人像与文本 prompt 共同决定主体和场景}{00:44:12}

右上参考图提供身份信息，主视频需要在新姿态、光照和场景中保持身份。简单复制脸部纹理会导致运动僵硬；过度服从文本又会丢失身份。应同时测 identity similarity、motion quality 与 prompt adherence。

\lecturefigure{slide-419-synchronized-audio.jpg}{Movie Gen Audio 根据视频与文本生成同步雷声和配乐}{00:44:43}

音频模型需要知道何时发生雷电、场景情绪以及哪些声音属于画内（diegetic）或画外。讲座明确指出，前面展示的 T2V 模型本身输出静音视频；同步音频来自独立的 Movie Gen Audio 模型。

\begin{importantbox}{Movie Gen 是模型家族，不是一张网络}
基础 T2V、video editing、personalization 和 audio generation 共享研究路线，但输入、目标、数据和后训练不同。评估或部署时必须说明使用哪个模型与哪条 pipeline，不能把家族能力全部归给 30B T2V backbone。
\end{importantbox}

\teachervoice{01:11:28--01:13:20，老师最后再次澄清：课堂主模型生成静音视频，音频来自独立模型。联合音视频生成理论上可共享信息，但高质量对齐数据比单独视频更稀缺。}

\subsection{本章小结}

Movie Gen 家族覆盖纯生成、编辑、个性化和音频，但这些能力不来自同一训练目标。基础 T2V 展示强视觉规律拟合；编辑强调内容保持；个性化强调身份；音频强调跨模态同步。能力命名必须跟具体模型和数据路径绑定。

